% from ega2-1-fr.tex
\oldpage[II]{167}
Il est immédiat (\hyperref[I.3.4.3-fr]{I, 3.4.3}) que la donnée d'une structure
de $Y$-schéma en groupes sur un $Y$-schéma $G$ équivaut à la donnée, pour
\emph{tout} $Y$-préschéma $Z$, d'une structure de \emph{groupe} sur
l'ensemble $\operatorname{Hom}_Y(Z,G)$, ces structures devant être telles que
pour tout $Y$-morphisme $Z\to Z'$, l'application correspondante
$\operatorname{Hom}_Y(Z',G)\to\operatorname{Hom}_Y(Z,G)$ soit un
homomorphisme de groupes. Dans le cas particulier de $G_m$ que nous
considérons ici, $\operatorname{Hom}_Y(Z,G)$ s'identifie à l'ensemble des
$Z$-sections de $Z\times_YG_m$ (\hyperref[I.3.3.14-fr]{I, 3.3.14}), donc à
l'ensemble des $Z$-sections de
$\operatorname{Spec}(\mathcal{O}_Z[T,T^{-1}])$~; finalement, le même
raisonnement que dans (\hyperref[I.3.3.15-fr]{I, 3.3.15}) montre que cet
ensemble s'identifie canoniquement à l'ensemble des éléments \emph{inversibles}
de l'anneau $\Gamma(Z,\mathcal{O}_Z)$, et la structure de groupe sur cet
ensemble est la structure provenant de la multiplication dans l'anneau
$\Gamma(Z,\mathcal{O}_Z)$. Le lecteur pourra vérifier que les morphismes $p$
et $s$ considérés plus haut sont obtenus de la façon suivante~: ils
correspondent d'après (\hyperref[II.1.2.7-fr]{1.2.7}) et
(\hyperref[II.1.4.6-fr]{1.4.6}) à des homomorphismes de $\mathcal{O}_Y$-Algèbres
\begin{align*}
  \pi &: \mathcal{O}_Y[T,T^{-1}]
  \to\mathcal{O}_Y[T,T^{-1},T',T'^{-1}],\\
  \sigma &: \mathcal{O}_Y[T,T^{-1}]\to\mathcal{O}_Y[T,T^{-1}]
\end{align*}
et sont entièrement définis par les données de $\pi(T)=TT'$ et
$\sigma(T)=T^{-1}$.

Cela étant, $G_m$ peut être considéré comme un «~\emph{domaine d'opérateurs
universel}~» pour tout \emph{cône affine}
$C=\operatorname{Spec}(\mathcal{S})$, où $\mathcal{S}$ est une
$\mathcal{O}_Y$-Algèbre graduée quasi-cohérente à degrés positifs. Cela
signifie qu'on peut définir canoniquement un $Y$-morphisme
$G_m\times_YC\to C$ qui a les propriétés formelles d'une loi externe d'un
ensemble muni d'un groupe d'opérateurs~; ou encore, comme ci-dessus pour les
schémas en groupes, on se donne pour tout $Y$-préschéma $Z$ une loi externe
sur $\operatorname{Hom}_Y(Z,C)$, ayant le groupe
$\operatorname{Hom}_Y(Z,G_m)$ comme ensemble d'opérateurs, avec les axiomes
usuels des ensembles munis d'un groupe d'opérateurs, et une condition de
compatibilité avec les $Y$-morphismes $Z\to Z'$. Dans le cas présent, le
morphisme $G_m\times_YC\to C$ est défini par la donnée de l'homomorphisme de
$\mathcal{O}_Y$-Algèbres
\[
  \mathcal{S}\to
  \mathcal{S}\otimes_{\mathcal{O}_Y}\mathcal{O}_Y[T,T^{-1}]
  =\mathcal{S}[T,T^{-1}]
\]
qui, à toute section $s_n\in\Gamma(U,\mathcal{S}_n)$ ($U$ ouvert de $Y$) fait
correspondre la section
$s_nT^n\in\Gamma(U,\mathcal{S}\otimes_{\mathcal{O}_Y}
\mathcal{O}_Y[T,T^{-1}])$.

Inversement, supposons donnée une $\mathcal{O}_Y$-Algèbre quasi-cohérente
$\mathcal{S}$ \emph{non graduée a priori}, et, sur
$C=\operatorname{Spec}(\mathcal{S})$, une structure de «~$Y$-schéma en
ensembles \emph{munis de groupes d'opérateurs}~» ayant pour domaine
d'opérateurs le $Y$-schéma en groupes $G_m$~; alors on en déduit
canoniquement une \emph{graduation} de $\mathcal{O}_Y$-Algèbre sur
$\mathcal{S}$. En effet, la donnée d'un $Y$-morphisme $G_m\times_YC\to C$
équivaut à celle d'un homomorphisme de $\mathcal{O}_Y$-Algèbres
$\psi:\mathcal{S}\to\mathcal{S}[T,T^{-1}]$, qui s'écrira donc
$\psi=\sum_{n\in\mathbf{Z}}\psi_nT^n$, où les
$\psi_n:\mathcal{S}\to\mathcal{S}$ sont des homomorphismes de
$\mathcal{O}_Y$-Modules (avec $\psi_n(s)=0$ sauf pour un nombre fini d'indices
pour toute section $s\in\Gamma(U,\mathcal{S})$, $U$ ouvert dans $Y$). On peut
alors vérifier que les axiomes des ensembles munis d'un groupe d'opérateurs
entraînent que les $\psi_n(\mathcal{S})=\mathcal{S}_n$ définissent une
graduation (à degrés positifs ou négatifs) de $\mathcal{O}_Y$-Algèbre sur
$\mathcal{S}$, les $\psi_n$ étant les projecteurs correspondants. On a ainsi
une notion de structure de «~\emph{cône affine}~» sur tout $Y$-schéma affine,
définie de façon «~géométrique~» sans appel à une graduation préalable.

